\documentclass{article}
\usepackage[T2A]{fontenc}
\usepackage[utf8]{inputenc}
\usepackage[russian]{babel}
\usepackage[margin=1cm]{geometry}
\usepackage{array}

\pagestyle{empty}
\setlength{\parindent}{0pt}
\setlength{\tabcolsep}{0pt}
\renewcommand{\arraystretch}{0.95}

% Четыре уникальных варианта
\newcommand{\taskA}{%
  Разработайте программу для ввода набора названий учебных предметов
  (проверяйте, что названия~--- непустые строки) и вывода тех из них,
  у которых длина названия превышает 10 символов.
  Используйте принципы S и D, а также свойства ООП.
  Для хранилища названий разработайте качественный набор unit-тестов.%
}
\newcommand{\taskB}{%
  Разработайте программу для ввода набора результатов измерения высоты тела
  (проверяйте, что высота~--- положительное число) и вывода тех из них,
  которые больше среднего арифметического.
  Используйте принципы S и D, а также свойства ООП.
  Для хранилища измерений разработайте качественный набор unit-тестов.%
}
\newcommand{\taskC}{%
  Разработайте программу для ввода набора телефонов абонентов
  (проверяйте, что каждый символ~--- цифра или пробел) и вывода тех из них,
  в которых количество цифр не больше 10.
  Используйте принципы S и D, а также свойства ООП.
  Для хранилища телефонов разработайте качественный набор unit-тестов.%
}
\newcommand{\taskD}{%
  Разработайте программу для ввода набора результатов бросания случайной точки
  на отрезок $[0;1]$ (проверяйте, что каждое число принадлежит этому отрезку)
  и вывода тех из них, для которых модуль разности текущего числа и $0{,}5$
  не превышает $0{,}1$.
  Используйте принципы S и D, а также свойства ООП.
  Для хранилища результатов разработайте качественный набор unit-тестов.%
}

\begin{document}

\small
\noindent
\begin{tabular}{@{}p{0.48\textwidth}@{\hspace{0.02\textwidth}}p{0.48\textwidth}@{}}
\textbf{Вариант 1.} \taskA & \textbf{Вариант 2.} \taskB \\[1mm] \hline
\textbf{Вариант 3.} \taskC & \textbf{Вариант 4.} \taskD \\[1mm] \hline
\textbf{Вариант 1.} \taskA & \textbf{Вариант 2.} \taskB \\[1mm] \hline
\textbf{Вариант 3.} \taskC & \textbf{Вариант 4.} \taskD \\[1mm] \hline
\textbf{Вариант 1.} \taskA & \textbf{Вариант 2.} \taskB \\[1mm] \hline
\textbf{Вариант 3.} \taskC & \textbf{Вариант 4.} \taskD \\[1mm] \hline
\textbf{Вариант 1.} \taskA & \textbf{Вариант 2.} \taskB \\[1mm] \hline
\textbf{Вариант 3.} \taskC & \textbf{Вариант 4.} \taskD \\[1mm] \hline
\textbf{Вариант 1.} \taskA & \textbf{Вариант 2.} \taskB \\[1mm] \hline
\textbf{Вариант 3.} \taskC & \textbf{Вариант 4.} \taskD \\[1mm] \hline
\end{tabular}

\end{document}